\documentclass[12pt,a4paper]{article}

\usepackage[spanish,es-tabla]{babel}
\usepackage[T1]{fontenc}
\usepackage[utf8]{inputenc}
\usepackage{lmodern}
\usepackage[a4paper,margin=2.5cm]{geometry}
\usepackage{setspace}
\usepackage{parskip}
\usepackage{hyperref}
\usepackage{graphicx}
\usepackage{float}
\usepackage{enumitem}
\usepackage{xcolor}

\hypersetup{
    colorlinks=true,
    linkcolor=blue,
    urlcolor=blue,
    pdftitle={Manual de Usuario Global de Instalaciones},
    pdfauthor={OpenAI Codex}
}

\setstretch{1.1}
\setlist[itemize]{topsep=4pt,itemsep=2pt}
\setlist[enumerate]{topsep=4pt,itemsep=2pt}

\newcommand{\placeholderimage}[2]{
    \begin{figure}[H]
        \centering
        \fbox{\parbox[c][6cm][c]{0.88\textwidth}{
            \centering
            \textbf{#1}\\[0.4cm]
            #2
        }}
        \caption{#2}
    \end{figure}
}

\newcommand{\manualimage}[2]{
    \IfFileExists{#1}{
        \begin{figure}[H]
            \centering
            \includegraphics[width=0.9\textwidth]{#1}
            \caption{#2}
        \end{figure}
    }{
        \placeholderimage{#1}{#2}
    }
}

\newcommand{\manualtripleimage}[4]{
    \IfFileExists{#1}{
        \begin{figure}[H]
            \centering
            \includegraphics[width=0.31\textwidth]{#1}\hfill
            \includegraphics[width=0.31\textwidth]{#2}\hfill
            \includegraphics[width=0.31\textwidth]{#3}
            \caption{#4}
        \end{figure}
    }{
        \placeholderimage{#1 / #2 / #3}{#4}
    }
}

\title{Manual de Usuario Global de Instalaciones}
\author{}
\date{\today}

\begin{document}

\maketitle
\tableofcontents
\newpage

\section{Finalidad del manual}

Este manual explica cómo funciona una instalación dentro de Allplan desde el punto de vista del usuario.

Se ha tomado la instalación de Agua como ejemplo, pero el funcionamiento general es común al resto de instalaciones. Lo que cambia entre unas y otras son algunos elementos concretos, ciertos nombres de layer y algunos atributos, pero la forma de trabajar es prácticamente la misma.

El objetivo de este documento es que cualquier usuario pueda entender:

\begin{itemize}
\item cómo empezar una instalación;
\item cómo dibujarla;
\item cómo editarla;
\item cómo aplicar layers y atributos;
\item cómo insertar elementos definidos, macros y soportes;
\item cómo funcionan las copias o duplicados;
\item y qué comprobaciones conviene hacer antes de finalizar.
\end{itemize}

\section{A quién va dirigido}

Este manual está pensado para usuarios que necesiten modelar instalaciones en Allplan y entender el comportamiento general de la herramienta, sin entrar en detalles de programación ni en nombres internos del sistema.

\section{Cómo se complementa este manual}

Este documento se organiza en dos grandes bloques:

\begin{enumerate}
\item El funcionamiento global de las instalaciones.
\item Las particularidades de cada instalación.
\end{enumerate}

En la primera parte se explica la lógica de uso común a todos los sistemas.

En la segunda parte se recogen las diferencias específicas de cada instalación, por ejemplo:

\begin{itemize}
\item los tipos de instalación disponibles;
\item las distribuciones admitidas;
\item los diámetros disponibles;
\item los accesorios automáticos;
\item los elementos definidos propios;
\item las limitaciones o avisos especiales;
\item y cualquier diferencia importante respecto a otras instalaciones.
\end{itemize}

De este modo, el usuario puede entender primero el funcionamiento general y después consultar, dentro del mismo manual, la unidad concreta de la instalación con la que vaya a trabajar.

\section{Qué partes son comunes en todas las instalaciones}

Todas las instalaciones comparten la misma lógica de trabajo:

\begin{enumerate}
\item Se selecciona el tipo de instalación.
\item Se configuran los parámetros principales en la paleta.
\item Se dibuja o edita el recorrido.
\item Se revisa la previsualización.
\item Se añaden, si hace falta, elementos definidos, macros o soportes.
\item Se aplican layers y atributos.
\item Se finaliza la creación.
\end{enumerate}

Lo que cambia entre instalaciones suele ser:

\begin{itemize}
\item el tipo de elementos disponibles;
\item los diámetros;
\item las conexiones;
\item algunos layers;
\item y algunos atributos asociados a cada sistema.
\end{itemize}

\section{Vista general del trabajo con una instalación}

El trabajo habitual con una instalación se basa en tres ideas:

\begin{itemize}
\item dibujar un recorrido;
\item enriquecer ese recorrido con información;
\item y generar el resultado final.
\end{itemize}

En la práctica, esto significa que primero se crea el trazado principal y después se añaden o ajustan el resto de elementos necesarios.

\manualimage{CAPTURA-01.png}{Paleta principal completa de la instalación con sus apartados visibles}

\section{Flujo de trabajo recomendado}

El orden más recomendable para trabajar es este:

\begin{enumerate}
\item Elegir el tipo de instalación.
\item Elegir el tipo de tubo o sistema.
\item Definir diámetro, distribución y resto de parámetros básicos.
\item Dibujar el recorrido principal.
\item Revisar uniones, codos y bifurcaciones.
\item Aplicar layers y atributos si es necesario.
\item Insertar elementos definidos, macros o soportes.
\item Comprobar que la información esté completa.
\item Pulsar el botón de finalizar.
\end{enumerate}

Seguir este orden ayuda a evitar correcciones posteriores innecesarias.

\section{Modos principales de trabajo}

La herramienta suele trabajar con tres modos principales.

\subsection{Modo creación}

Es el modo que se utiliza para dibujar un recorrido nuevo.

En este modo el usuario va definiendo puntos y la instalación crea el trazado entre ellos.

\subsection{Modo edición}

Es el modo que se utiliza para corregir o modificar un recorrido ya dibujado.

En este modo se pueden seleccionar tramos, borrar partes, cambiar diámetros o ajustar la geometría.

\subsection{Modo extender}

Es el modo que se utiliza para continuar un recorrido ya existente sin empezar desde cero.

\section{Qué debe configurar el usuario antes de dibujar}

Antes de empezar a dibujar conviene revisar los campos principales de la paleta.

\subsection{Tipo de instalación}

Es el campo donde el usuario elige el sistema con el que va a trabajar.

En Agua, por ejemplo, puede elegir entre varias opciones como:

\begin{itemize}
\item Polietilè;
\item Multicapa;
\item Armaflex.
\end{itemize}

Este campo define qué geometría y qué comportamiento tendrá la instalación.

\manualimage{CAPTURA-02.png}{Desplegable del tipo de instalación abierto con las opciones visibles}

\subsection{Diámetro}

El diámetro condiciona:

\begin{itemize}
\item el tamaño del elemento principal;
\item las piezas de conexión compatibles;
\item y parte de los atributos asociados.
\end{itemize}

Antes de dibujar, conviene comprobar que el diámetro seleccionado es el correcto.

\subsection{Distribución}

La distribución define cómo debe comportarse la instalación dentro del sistema.

En Agua es habitual trabajar con modos como:

\begin{itemize}
\item TD;
\item IS.
\end{itemize}

La distribución puede afectar a:

\begin{itemize}
\item la forma de los elementos;
\item los atributos aplicados;
\item y el comportamiento de ciertos accesorios.
\end{itemize}

\subsection{Tipo de agua u otras variantes}

Algunas instalaciones incluyen campos adicionales para distinguir variantes internas del sistema. Si la paleta muestra este dato, conviene definirlo antes de empezar.

\subsection{Ángulos y orientación}

Si la instalación trabaja con limitación angular o con orientación definida, estos parámetros deben comprobarse al inicio para evitar tener que rehacer el recorrido más adelante.

\section{Cómo dibujar una instalación}

El dibujo de la instalación se basa en una secuencia de clics que definen el recorrido.

\subsection{Inicio del recorrido}

El usuario hace clic en el primer punto del trazado.

\subsection{Continuación del recorrido}

Cada clic adicional añade un nuevo tramo.

La herramienta va construyendo la instalación entre esos puntos y mostrando una previsualización del resultado.

\subsection{Qué ve el usuario mientras dibuja}

Mientras se dibuja, la herramienta no solo muestra la línea del recorrido, sino también el comportamiento previsto de la instalación:

\begin{itemize}
\item tramos principales;
\item cambios de dirección;
\item uniones;
\item bifurcaciones.
\end{itemize}

\manualimage{CAPTURA-03.png}{Ejemplo de instalación en fase de dibujo con previsualización activa}

\section{Cómo se generan automáticamente los accesorios}

Una vez definido el recorrido, la herramienta interpreta la geometría y coloca automáticamente los accesorios necesarios.

\subsection{Codos}

Cuando el recorrido cambia de dirección, la herramienta genera el codo correspondiente.

\subsection{Manguitos o uniones}

Cuando dos tramos deben unirse de forma alineada, la herramienta genera la unión correspondiente.

\subsection{Bifurcaciones o tes}

Cuando un punto conecta tres direcciones, la herramienta genera la bifurcación adecuada.

\subsection{Regla general de prioridad}

Si un punto funciona como bifurcación, prevalece la bifurcación sobre otras soluciones.

Esto significa que el sistema no coloca varias piezas incompatibles en el mismo punto.

\manualtripleimage{CAPTURA-04A.png}{CAPTURA-04B.png}{CAPTURA-04C.png}{Ejemplo comparativo de un codo, una unión y una bifurcación}

\section{Cómo editar una instalación ya dibujada}

Una vez creado el recorrido, el usuario puede modificarlo.

\subsection{Seleccionar tramos}

Se pueden seleccionar:

\begin{itemize}
\item tramos individuales;
\item varios tramos a la vez;
\item o zonas concretas del recorrido.
\end{itemize}

\subsection{Borrar una sección}

El botón de borrar elimina el tramo o la parte seleccionada.

Después del borrado, la herramienta vuelve a analizar la conexión entre los tramos que quedan y reajusta automáticamente las piezas necesarias.

\subsection{Cambiar el diámetro}

El botón de modificar diámetro actualiza el diámetro del tramo o de la selección activa.

Después del cambio, la instalación vuelve a reconstruir el resultado para adaptarlo al nuevo valor.

\section{Layers}

El usuario puede trabajar con los layers definidos por defecto o aplicar otros manualmente.

\subsection{Layers por defecto}

Cada instalación ya trae una configuración base de layers.

\subsection{Aplicar layers manualmente}

Desde la paleta, el usuario puede seleccionar un layer y aplicar ese layer a los elementos seleccionados.

Esto es útil cuando se necesita clasificar el modelo de una forma concreta antes de finalizar.

\manualimage{CAPTURA-05.png}{Ejemplo del apartado de layers en la paleta y su aplicación sobre la instalación}

\section{Atributos}

Además de la geometría, la instalación puede llevar información adicional en forma de atributos.

\subsection{Atributos automáticos}

Parte de los atributos se asignan automáticamente según el tipo de instalación, el elemento y la configuración elegida.

\subsection{Atributos aplicados por el usuario}

La paleta también permite aplicar atributos manualmente a una selección.

Esto es útil cuando el usuario necesita completar o corregir información antes de crear el resultado final.

\subsection{Atributo padre}

En determinadas instalaciones existe un atributo principal que sirve para organizar y relacionar elementos. Este dato es importante porque también puede influir en las copias o duplicados automáticos.

\section{Qué revisar antes de finalizar}

Antes de pulsar el botón de finalizar, se recomienda comprobar:

\begin{itemize}
\item que el recorrido es correcto;
\item que los diámetros son los adecuados;
\item que los layers necesarios están aplicados;
\item que los atributos importantes están informados;
\item y que los elementos definidos están bien colocados.
\end{itemize}

Si falta información, la herramienta puede mostrar un aviso antes de continuar.

\manualimage{CAPTURA-06.png}{Mensaje de aviso previo a finalizar cuando falta información}

\section{Rotación y orientación}

Algunos elementos admiten rotación y orientación.

Esto se aplica especialmente a:

\begin{itemize}
\item elementos definidos;
\item macros;
\item ciertos accesorios que requieren una posición concreta.
\end{itemize}

Cuando la paleta muestra campos de rotación u orientación, el usuario puede utilizarlos para ajustar la posición final del elemento.

\manualimage{CAPTURA-07.png}{Ejemplo de orientación o rotación aplicada a un elemento}

\section{Macros}

La herramienta permite colocar macros de librería dentro de la instalación.

\subsection{Para qué sirven}

Las macros permiten insertar objetos ya preparados, por ejemplo elementos de librería que deben colocarse en puntos concretos del recorrido.

\subsection{Qué debe hacer el usuario}

El flujo habitual es:

\begin{enumerate}
\item Elegir el tipo de macro.
\item Elegir el archivo correspondiente.
\item Elegir dónde se colocará.
\item Ajustar altura y rotación si es necesario.
\item Insertar la macro.
\end{enumerate}

\subsection{Dónde puede colocarse}

La macro puede situarse:

\begin{itemize}
\item al inicio del recorrido;
\item al final;
\item o en un punto libre.
\end{itemize}

\subsection{Altura de la macro}

La altura puede definirse directamente o en relación con un local seleccionado, según el caso.

\manualimage{CAPTURA-08.png}{Bloque de la paleta para insertar macros}

\manualimage{CAPTURA-09.png}{Ejemplo de macro previsualizada dentro de la instalación}

\section{Elementos definidos}

Además del recorrido principal, la herramienta permite añadir elementos definidos propios de la instalación.

En Agua, por ejemplo, pueden existir elementos como válvulas, salidas o piezas concretas del sistema.

\subsection{Cómo se insertan}

Hay dos formas habituales:

\begin{itemize}
\item insertarlos sobre un punto del recorrido;
\item o insertarlos como punto libre.
\end{itemize}

\subsection{Qué debe hacer el usuario}

El usuario debe:

\begin{enumerate}
\item elegir el tipo de elemento;
\item seleccionar el punto de inserción;
\item ajustar rotación o altura si hace falta;
\item confirmar la colocación.
\end{enumerate}

\manualimage{CAPTURA-10.png}{Bloque de la paleta para insertar Elementos Definidos}

\manualimage{CAPTURA-11.png}{Inserción de un elemento definido como punto libre}

\section{Puntos no definidos}

Los puntos no definidos sirven para marcar la lógica del recorrido sin colocar todavía un elemento final.

Son útiles para describir:

\begin{itemize}
\item inicios;
\item finales;
\item intermedios;
\item bifurcaciones;
\item conexiones entre caminos.
\end{itemize}

\subsection{Para qué sirven}

Ayudan a organizar recorridos complejos y a definir la intención del trazado antes de convertirlo en una solución completa.

\subsection{Puntos comunes}

Si varios caminos coinciden en un punto o pasan muy cerca, la herramienta puede interpretar que existe una relación entre ellos.

Esto ayuda a construir una topología coherente.

\manualimage{CAPTURA-12.png}{Ejemplo de puntos no definidos en distintos colores}

\manualimage{CAPTURA-13.png}{Ejemplo de punto común entre dos caminos}

\section{Soportes}

La herramienta también permite trabajar con soportes.

\subsection{Qué hace el usuario}

El flujo habitual es:

\begin{enumerate}
\item configurar el soporte en la paleta;
\item elegir el tipo de soporte y revisar sus parámetros;
\item pulsar el botón para insertarlo;
\item definir su posición con los puntos necesarios;
\item revisar la previsualización;
\item acumularlo o crearlo.
\end{enumerate}

Para que el soporte quede bien definido, conviene completar antes de insertarlo los campos de tipo, superficie, cotas y ángulo.

\subsection{Tipos de soporte}

De forma general, la herramienta permite trabajar con familias de soporte como:

\begin{itemize}
\item \texttt{Zeta}
\item \texttt{Omega}
\item \texttt{Cinta}
\end{itemize}

Cada una responde a una solución física distinta. Por eso, antes de colocarlo, el usuario debe elegir el tipo que corresponda al sistema y al montaje que quiere representar.

En función de la instalación, pueden aparecer variantes o subtipos asociados a cada familia.

\subsection{Subtipo de instalación}

Además del tipo de soporte, la paleta puede mostrar el campo \texttt{Subtipo instalación}.

Este campo ayuda a identificar con qué familia de instalación se relaciona el soporte. Es útil para mantener coherencia entre el soporte que se inserta y la instalación sobre la que se está trabajando.

\manualimage{CAPTURA-14B.png}{Tipo de soporte, subtipo instalación y superficie en la paleta de soportes}

\subsection{Superficie}

El campo \texttt{Superficie} permite definir cómo debe considerarse el soporte:

\begin{itemize}
\item \texttt{Liso}
\item \texttt{Perforado}
\end{itemize}

Esta elección forma parte de la definición final del soporte y de la información que queda asociada a él. Por eso conviene revisarla antes de insertarlo.

No todos los tipos de soporte utilizan esta distinción exactamente de la misma manera, pero para el usuario la regla práctica es sencilla: debe escoger la superficie que corresponda al soporte real que quiere representar.

\subsection{Cota A y Cota B}

Estas dos medidas son básicas para definir la geometría del soporte:

\begin{itemize}
\item \texttt{Cota A}: indica la altura del soporte o su desarrollo vertical.
\item \texttt{Cota B}: indica el largo del soporte o su desarrollo horizontal.
\end{itemize}

En soportes con parte vertical, \texttt{Cota A} tiene un papel especialmente importante. En otros tipos, como ciertas cintas, no siempre interviene de la misma forma.

\subsection{Ángulo de inclinación}

El campo \texttt{Ángulo de inclinación} permite girar o inclinar el soporte respecto a su posición base.

Es útil cuando el soporte no debe quedar en una orientación neutra y necesita un ajuste adicional para adaptarse al montaje real.

Lo recomendable es definir primero la posición del soporte y usar después el ángulo como ajuste fino.

\manualimage{CAPTURA-14C.png}{Ejemplo de Cota A, Cota B y ángulo de inclinación en la paleta}

\subsection{Modos de edición}

Una vez que hay soportes acumulados, la herramienta permite trabajar con distintos modos de edición:

\begin{itemize}
\item \texttt{Desactivado}: la herramienta no entra en edición de soportes y los clics siguen el flujo normal de trabajo.
\item \texttt{Edición}: permite seleccionar soportes para revisarlos, borrarlos o aplicarles atributos.
\item \texttt{Edición Mover}: permite recolocar un soporte ya insertado o acumulado.
\end{itemize}

Estos modos son útiles cuando el recorrido principal ya está dibujado y solo falta ajustar la posición o la información de los soportes.

\manualimage{CAPTURA-14D.png}{Modos de edición de soportes en la paleta}

\subsection{Atributos de soporte}

La paleta también incluye un campo para aplicar el \texttt{Atributo de soporte}.

Este atributo sirve para identificar, clasificar o completar la información de los soportes antes de la creación final.

Si el usuario tiene soportes seleccionados, el atributo se aplica a esa selección. Si no hay una selección activa, la aplicación afecta a los soportes acumulados dentro de la operación en curso.

Esto resulta especialmente útil cuando se quiere dejar todos los soportes correctamente preparados antes de finalizar la instalación.

\manualimage{CAPTURA-14E.png}{Campo de atributo de soporte y aplicación sobre soportes seleccionados o acumulados}

\subsection{Cuándo conviene colocarlos}

Normalmente es más cómodo colocar los soportes cuando el recorrido principal ya está bastante definido.

También conviene tener en cuenta que, antes de acumular o crear un soporte, primero debe haberse posicionado correctamente. Si todavía no se ha definido su colocación, la herramienta avisa y no lo añade.

\manualimage{CAPTURA-14.png}{Bloque de soportes en la paleta y preview de un soporte}

\section{Cómo funcionan los duplicados o copias}

Este punto es importante porque puede generar dudas.

\subsection{Duplicados no deseados}

La herramienta evita repetir elementos iguales cuando no corresponde.

Esto ayuda a que no aparezcan piezas duplicadas por error en el resultado final.

\subsection{Copias intencionadas}

En algunos casos, la instalación puede generar copias en otros archivos de dibujo según la información asociada al elemento.

Estas copias no son un error. Forman parte del funcionamiento previsto de la herramienta.

De forma general, el comportamiento es este:

\begin{enumerate}
\item el elemento se crea en el archivo principal en el que el usuario está trabajando;
\item si ese elemento tiene informado un atributo padre, la herramienta puede buscar un archivo de dibujo relacionado con esa referencia;
\item si el archivo de destino está cargado y coincide con esa referencia, la herramienta crea allí una copia relacionada;
\item esa copia queda vinculada a la actualización posterior de la instalación.
\end{enumerate}

\subsection{Qué debe entender el usuario}

El usuario debe saber que:

\begin{itemize}
\item puede existir un elemento en el archivo principal;
\item y puede existir una copia relacionada en otro archivo de dibujo.
\end{itemize}

Para que esa copia pueda producirse, normalmente deben cumplirse estas condiciones:

\begin{itemize}
\item el atributo padre debe estar informado;
\item el archivo de dibujo de destino debe estar cargado;
\item y el nombre del archivo de destino debe coincidir con la referencia que utiliza la instalación.
\end{itemize}

Si el archivo de destino no está cargado o no coincide con esa referencia, la copia no se genera.

\subsection{Qué pasa cuando la instalación se vuelve a editar}

Cuando la instalación se actualiza, esas copias también se actualizan para que no queden versiones antiguas.

Antes de crear las nuevas copias, la herramienta localiza las anteriores, las elimina de los archivos correspondientes y vuelve a generar las copias actualizadas.

\manualimage{CAPTURA-15.png}{Ejemplo de un elemento original y su copia relacionada en otro archivo}

\section{Qué ocurre al pulsar el botón de finalizar}

Cuando el usuario pulsa el botón de finalizar, la herramienta:

\begin{enumerate}
\item revisa la información disponible;
\item reconstruye los elementos necesarios;
\item aplica atributos y layers;
\item incorpora elementos definidos, macros o soportes;
\item crea el resultado definitivo;
\item y actualiza las copias relacionadas si las hubiera.
\end{enumerate}

En ese momento la instalación pasa de estar en fase de preparación a quedar creada como resultado final en el documento.

\section{Recomendaciones prácticas de uso}

\begin{itemize}
\item Configurar primero el sistema antes de empezar a dibujar.
\item No cambiar de criterio de diámetro continuamente durante el dibujo si no es necesario.
\item Revisar la distribución antes de finalizar.
\item Aplicar layers y atributos cuando el recorrido ya esté claro.
\item Insertar macros y elementos definidos cuando la base principal esté estable.
\item Revisar los avisos antes de aceptar la creación final.
\end{itemize}

\section{Resumen final}

La herramienta de instalaciones está pensada para que el usuario dibuje un recorrido, lo complete con la información necesaria y genere un resultado final coherente dentro de Allplan.

Aunque cada instalación tenga sus particularidades, el funcionamiento general es siempre el mismo:

\begin{itemize}
\item configurar;
\item dibujar;
\item revisar;
\item completar;
\item y finalizar.
\end{itemize}

Por eso, entendiendo bien el ejemplo de Agua, se entiende también el comportamiento general del resto de instalaciones.

\section{Particularidades de cada instalación}

Una vez entendido el funcionamiento general, es importante conocer las diferencias propias de cada instalación.

Estas diferencias no cambian la forma general de trabajar, pero sí modifican:

\begin{itemize}
\item los sistemas disponibles;
\item los diámetros;
\item los accesorios automáticos;
\item los elementos definidos;
\item y algunas reglas concretas de uso.
\end{itemize}

\section{Instalación de Agua}

\noindent En esta unidad solo se recogen las particularidades de Agua. El uso general de la herramienta ya se ha explicado en los apartados anteriores, por lo que aquí se detallan únicamente las reglas, avisos y casos propios de esta instalación.

\subsection{Sistemas disponibles en Agua}

En la instalación de Agua, el usuario puede trabajar con estos sistemas:

\begin{itemize}
\item Polietilè;
\item Multicapa;
\item Armaflex.
\end{itemize}

Cada uno de ellos pertenece a la instalación de Agua, pero representa una solución distinta dentro del modelo.

\manualimage{CAPTURA-AGUA-01.png}{Desplegable con los sistemas disponibles en Agua}

\subsection{Distribuciones en Agua}

En Agua el usuario puede trabajar con estas distribuciones:

\begin{itemize}
\item EN;
\item TD;
\item IS.
\end{itemize}

La distribución condiciona el comportamiento del sistema, los layers aplicados y el catálogo de piezas que puede resolverse automáticamente.

Cuando se selecciona la distribución \texttt{EN}, la paleta muestra además el campo \texttt{Cara (EN)}, que permite indicar si el caso corresponde a \texttt{Cara X} o \texttt{Cara Y}.

\subsubsection{Combinaciones que requieren atención}

\begin{itemize}
\item Conviene elegir la distribución antes de empezar a dibujar.
\item Si la distribución se cambia después, es necesario revisar de nuevo codos, uniones, bifurcaciones y layers.
\item No todos los sistemas de Agua están disponibles en todas las distribuciones.
\end{itemize}

En la configuración actual:

\begin{itemize}
\item \texttt{Multicapa} no está disponible en \texttt{IS}.
\item \texttt{Armaflex} no está disponible en \texttt{IS}.
\end{itemize}

Si el usuario intenta trabajar con alguna de esas combinaciones, la herramienta muestra un aviso indicando que ese tipo de tubo no existe o no está disponible para la distribución \texttt{IS}, y pide volver a \texttt{Modo creación}, cambiar la distribución a \texttt{TD} y regresar a \texttt{Modo configuración}.

\manualimage{CAPTURA-AGUA-02.png}{Campo de distribución en Agua con la opción EN y el campo Cara (EN) visible}

\manualimage{CAPTURA-AGUA-03.png}{Aviso mostrado al intentar usar Multicapa o Armaflex en distribución IS}

\subsection{Diámetros y cambios de diámetro en Agua}

En la instalación de Agua, los diámetros más habituales en la configuración actual son:

\begin{itemize}
\item 20;
\item 25.
\end{itemize}

El diámetro afecta a:

\begin{itemize}
\item el tamaño del tubo;
\item los accesorios compatibles;
\item la geometría final;
\item y parte de la información asociada al elemento.
\end{itemize}

\subsubsection{Qué ocurre al cambiar el diámetro de un tramo}

Cuando se modifica el diámetro de un único tramo, la herramienta muestra una ventana de confirmación con:

\begin{itemize}
\item la ruta;
\item el segmento;
\item el diámetro anterior;
\item y el nuevo diámetro.
\end{itemize}

El usuario debe confirmar el cambio antes de que la geometría se regenere.

\manualimage{CAPTURA-AGUA-04.png}{Confirmación de cambio de diámetro sobre un único tramo}

\subsubsection{Qué ocurre al cambiar el diámetro de varios tramos}

Si hay varios tramos seleccionados, la ventana de confirmación informa de:

\begin{itemize}
\item cuántos segmentos se van a modificar;
\item qué diámetro o diámetros tenían antes;
\item y qué nuevo diámetro se va a aplicar.
\end{itemize}

Después de aceptar, la instalación reconstruye el resultado con el nuevo valor.

\manualimage{CAPTURA-AGUA-05.png}{Confirmación de cambio de diámetro sobre varios tramos seleccionados}

\subsubsection{Qué ocurre si el cambio afecta a un codo}

Si el cambio de diámetro afecta a un giro resuelto con codo, no aparece un mensaje especial para el codo.

Lo que se muestra es el mensaje general de cambio de diámetro y, al aceptarlo, la herramienta recalcula el encuentro para regenerar el codo con la nueva condición.

\manualimage{CAPTURA-AGUA-06.png}{Ejemplo de codo recalculado después de modificar el diámetro del tramo}

\subsubsection{Qué ocurre si hay dos tramos consecutivos en la misma dirección con distinto diámetro}

Si dos segmentos consecutivos siguen en la misma dirección pero cambian de diámetro, la instalación no coloca un manguito normal, sino un \texttt{manguito reductor}.

Este caso es importante porque es la forma habitual de resolver una transición recta entre dos diámetros distintos.

Conviene revisarlo siempre en pantalla para comprobar que la reducción se ha generado justo en el punto esperado.

\manualimage{CAPTURA-AGUA-07.png}{Ejemplo de manguito reductor entre dos tramos consecutivos en la misma dirección}

\subsubsection{Qué ocurre si después de un codo el siguiente tramo tiene otro diámetro}

Si el recorrido llega a un \texttt{codo} con un diámetro y el tramo que sale de ese codo tiene otro distinto, la herramienta muestra igualmente el mensaje general de confirmación del cambio de diámetro.

Es decir, el usuario verá la misma confirmación de cambio de diámetro que en los demás casos, ya sea para un único tramo o para varios tramos seleccionados. No aparece un aviso exclusivo para el caso \texttt{codo + cambio de diámetro}.

Después de aceptar, la instalación mantiene la lógica del giro y recalcula el encuentro.

En la lógica interna actual de Agua, el codo se resuelve tomando como referencia el diámetro mayor de los dos lados del giro para ajustar ese encuentro.

Por eso, este caso conviene revisarlo siempre en pantalla antes de finalizar, especialmente si el usuario esperaba una advertencia específica o una transición distinta.

\manualimage{CAPTURA-AGUA-07B.png}{Mensaje de cambio de diámetro y resultado cuando tras un codo el siguiente tramo tiene otro diámetro}

\subsection{Accesorios automáticos en Agua}

En Agua los accesorios automáticos dependen de la geometría del recorrido y de los diámetros de los tramos que se encuentran.

\subsubsection{Codos}

En Agua la resolución automática de giros se hace con \texttt{codos de 90 grados}.

Esto significa que no debe esperarse un codo automático de \texttt{45 grados} dentro de esta instalación.

\manualimage{CAPTURA-AGUA-08.png}{Ejemplo de codo de 90 grados generado automáticamente en Agua}

\subsubsection{Manguitos}

Cuando dos tramos consecutivos están alineados y mantienen el mismo diámetro, la herramienta coloca un manguito o unión recta.

\subsubsection{Manguitos reductores}

Cuando esos dos tramos alineados tienen distinto diámetro, la unión se resuelve mediante un manguito reductor.

Cuando el cambio de diámetro se produce en un tramo recto, esta es la solución habitual. En cambio, si el cambio aparece justo a la salida de un codo, conviene revisar el resultado porque en ese caso manda la lógica del propio giro.

\subsubsection{Tes o bifurcaciones}

Cuando en un mismo punto confluyen tres direcciones, la instalación intenta resolver el encuentro mediante una \texttt{Te}.

En Agua conviene distinguir dos grandes familias:

\begin{itemize}
\item \texttt{Te iguales}, cuando las tres bocas trabajan con el mismo diámetro.
\item \texttt{Te reducidas}, cuando una de las bocas trabaja con un diámetro distinto.
\end{itemize}

En la configuración actual, los casos más representativos son:

\begin{itemize}
\item \texttt{Te Ø20};
\item \texttt{Te Ø25};
\item \texttt{Te Ø25-20-25};
\item \texttt{Te Ø25-20-20};
\item \texttt{Te Ø25-25-20}.
\end{itemize}

Si el nodo se resuelve como \texttt{Te}, en ese mismo punto no se coloca además un codo ni un manguito.

\manualimage{CAPTURA-AGUA-09.png}{Ejemplo comparativo de los distintos tipos de Te utilizados en Agua}

\subsubsection{Qué ocurre si no existe una Te compatible}

Si la combinación de diámetros de entrada, salida y rama no tiene una \texttt{Te} disponible, la herramienta muestra un aviso indicando que no existe una \texttt{TE} para la combinación de diámetros seleccionada y pide modificar los diámetros de los segmentos para que coincidan con un tipo disponible.

\manualimage{CAPTURA-AGUA-10.png}{Aviso mostrado cuando no existe una Te para la combinación de diámetros seleccionada}

\subsection{Longitudes mínimas y máximas de tubo en Agua}

Además del sistema, la distribución y el diámetro, en Agua también hay condiciones de longitud que el usuario debe conocer.

\subsubsection{Longitud mínima}

La longitud mínima de tramo en la configuración actual es \texttt{350 mm}.

Si el usuario dibuja un segmento más corto:

\begin{itemize}
\item la herramienta muestra el tipo de instalación;
\item indica la longitud mínima requerida;
\item indica la longitud real del tramo dibujado;
\item y pregunta si se desea ignorar la restricción o cancelar para recolocar el punto.
\end{itemize}

Esto permite decidir en el momento si se mantiene ese tramo corto o si se corrige antes de seguir dibujando.

\manualimage{CAPTURA-AGUA-11.png}{Aviso de longitud mínima no cumplida al dibujar un tramo demasiado corto}

\subsubsection{Longitud máxima}

La longitud máxima de referencia en Agua es \texttt{5,00 m}.

Si un tramo supera esa longitud, la herramienta no muestra un aviso bloqueante, sino que divide automáticamente ese tramo en subtramos más cortos para poder resolver la instalación dentro del límite permitido.

Por eso, cuando se trabaja con recorridos largos, conviene revisar el resultado para comprobar dónde se han producido los cortes y cómo han quedado las uniones generadas.

\manualimage{CAPTURA-AGUA-12.png}{Ejemplo de tramo largo dividido automáticamente al superar la longitud máxima}

\subsection{Layers y atributo padre en Agua}

En Agua hay dos temas que conviene entender juntos:

\begin{itemize}
\item los \texttt{layers} de la instalación;
\item el valor que se aplica como \texttt{atributo padre};
\end{itemize}

\subsubsection{Layers propios de Agua}

En la configuración actual, Agua trabaja con una base de layers como:

\begin{itemize}
\item \texttt{AIGUA}
\item \texttt{AIGUA CARA X}
\item \texttt{AIGUA CARA Y}
\item \texttt{IS CON AIGUA FAB}
\item \texttt{IS CON AIGUA OBR}
\end{itemize}

Además, la instalación parte de un layer por defecto para la geometría principal y de un layer específico para la polilínea de eje cuando esta se genera como apoyo.

Por eso, aunque el usuario vea solo un desplegable de layers en la paleta, detrás hay una lógica propia de Agua que conviene respetar.

\manualimage{CAPTURA-AGUA-17.png}{Desplegable de layers de Agua mostrando las opciones disponibles en paleta}

\subsubsection{Qué debe entender el usuario sobre los layers en Agua}

Si el usuario no aplica un layer manualmente, la instalación utiliza su configuración por defecto.

Si el usuario selecciona elementos y usa \texttt{Aplicar layer}, esa asignación queda guardada para esos elementos y se utiliza al crear el resultado final.

Esto es especialmente importante en Agua cuando:

\begin{itemize}
\item se quiere distinguir fabricación y obra;
\item se necesita separar \texttt{Cara X} y \texttt{Cara Y};
\item o se quiere corregir la clasificación de un tramo o accesorio antes de finalizar.
\end{itemize}

En general, después de cambiar distribución, geometría o selección de elementos, conviene revisar también los layers.

\subsubsection{Casos en los que conviene revisar el layer final}

En Agua hay piezas, sobre todo en algunos casos \texttt{TD}, en las que la geometría puede resolverse con varias partes y no siempre interesa forzar todas ellas con el mismo criterio visual.

Por eso, si después de aplicar un layer el usuario ve una pieza que no responde exactamente como esperaba, conviene revisar el resultado final en pantalla en lugar de asumir que se trata de un error.

La regla práctica para el manual es esta:

\begin{itemize}
\item el layer aplicado desde paleta manda en el flujo habitual;
\item pero algunas piezas técnicas pueden conservar parte de su lógica propia de modelo.
\end{itemize}

\manualimage{CAPTURA-AGUA-18.png}{Ejemplo de layer aplicado en Agua sobre tramos o accesorios seleccionados}

\subsubsection{Qué es el atributo padre en Agua}

En Agua, el \texttt{Valor atributo} que el usuario aplica desde la paleta no es solo una etiqueta descriptiva.

Ese valor se utiliza como \texttt{atributo padre}, que sirve para:

\begin{itemize}
\item identificar elementos que pertenecen a una misma referencia;
\item organizar la instalación;
\item y validar que la información esté completa antes de finalizar.
\end{itemize}

En Agua, la lógica es esta:

\begin{itemize}
\item en \texttt{TD}, el valor aplicado se utiliza como \texttt{pmp\_pare};
\item en \texttt{IS}, ese mismo valor se utiliza como \texttt{pmp\_pare} y también como \texttt{6\_CC\_IS}.
\end{itemize}

Es decir, el usuario escribe un único valor, pero la instalación lo reutiliza según la distribución activa.

\subsubsection{Cómo aplica el usuario el atributo padre}

El flujo práctico es:

\begin{enumerate}
\item seleccionar el tramo o los elementos que deban compartir referencia;
\item escribir el valor en el campo \texttt{Valor atributo};
\item pulsar \texttt{Aplicar atributo}.
\end{enumerate}

Cuando se hace esto, la herramienta confirma que los atributos han sido aplicados y guarda esa información para la creación final.

Si el usuario no asigna layer o atributo padre a ciertos elementos, antes de finalizar puede aparecer un aviso indicando que hay elementos sin configuración completa.

\manualimage{CAPTURA-AGUA-19.png}{Aplicación de Valor atributo en Agua y aviso de elementos sin layer o sin atributo padre}

La mecánica de copia a otros archivos es común al sistema y se explica en el apartado global \texttt{Cómo funcionan los duplicados o copias}.

\subsection{Elementos definidos propios de Agua}

Además del trazado automático, Agua permite insertar elementos definidos propios del sistema.

En la configuración actual, Agua trabaja con elementos definidos como:

\begin{itemize}
\item Colze Base;
\item Clau de Pas;
\item Taps;
\item T sortida.
\end{itemize}

Estos elementos se utilizan para completar la instalación con piezas concretas del sistema una vez que el recorrido principal ya está claro.

Además, algunos de estos elementos pueden mostrar un aviso si se intentan colocar en una distribución no admitida, especialmente en \texttt{IS}.

\manualimage{CAPTURA-AGUA-13.png}{Paleta de Elementos Definidos en Agua con los tipos disponibles}

\subsubsection{Tipo de punto en Elementos Definidos}

Además de elegir el elemento, en Agua es importante revisar el campo \texttt{Tipo de punto}, porque no todos los elementos pueden colocarse en cualquier posición de la instalación.

En la configuración actual, la relación es esta:

\begin{itemize}
\item \texttt{Colze base}: \texttt{Inicio} o \texttt{Final}.
\item \texttt{Taps}: \texttt{Inicio} o \texttt{Final}.
\item \texttt{T sortida}: \texttt{Intermedio ordenado} o \texttt{Intermedio libre}.
\item \texttt{Clau de Pas}: \texttt{Intermedio ordenado} o \texttt{Intermedio libre}.
\end{itemize}

\manualimage{CAPTURA-AGUA-13B.png}{Campo Tipo de punto en Elementos Definidos mostrando las opciones según el elemento seleccionado}

\subsubsection{Qué significa cada tipo de punto}

\begin{itemize}
\item \texttt{Inicio}: el elemento se coloca en el arranque de la polilínea o del recorrido correspondiente.
\item \texttt{Final}: el elemento se coloca en el extremo final del recorrido.
\item \texttt{Intermedio ordenado}: el elemento se coloca en una posición intermedia asociada al orden del recorrido.
\item \texttt{Intermedio libre}: el elemento se coloca en un punto intermedio elegido libremente por el usuario.
\end{itemize}

En otras palabras, \texttt{Inicio} y \texttt{Final} se usan para piezas que deben rematar o arrancar el recorrido, mientras que los tipos \texttt{Intermedio} se usan para piezas que deben quedar dentro del desarrollo de la instalación.

\subsubsection{Qué debe tener en cuenta el usuario}

\begin{itemize}
\item Si el elemento es de \texttt{Inicio/Final}, no debe intentarse como punto intermedio.
\item Si el elemento es intermedio, no debe intentarse como arranque o remate del recorrido.
\item Cuando la combinación no es válida, la herramienta avisa y no permite continuar con esa colocación.
\end{itemize}

Esto es especialmente importante en Agua porque el tipo de punto forma parte de la lógica de colocación del elemento, no es solo una etiqueta descriptiva.

\subsection{Soportes en Agua}

Además de las reglas comunes de soportes, en Agua hay particularidades de clasificación y de geometría que conviene revisar.

\subsubsection{Qué hace especial a los soportes de Agua}

En Agua no basta con fijarse en la forma del soporte. También hay que comprobar cómo queda asociado dentro de la instalación.

Esto es importante porque algunos soportes comparten familia geométrica con otras instalaciones. El caso más claro es \texttt{Omega}, ya que Agua y Saneamiento pueden trabajar con la misma base de soporte tipo \texttt{Varifix}.

Por eso, antes de insertar, conviene revisar que el campo \texttt{Subtipo instalación} quede realmente asociado a \texttt{Agua}. Si aparece otro subtipo, el soporte puede quedar clasificado como perteneciente a otra instalación aunque visualmente se parezca.

La confirmación previa a la inserción es un buen momento para comprobar \texttt{Tipo}, \texttt{Subtipo instalación}, \texttt{Superficie}, \texttt{Cota A}, \texttt{Cota B} y \texttt{Ángulo de inclinación} antes de continuar.

\manualimage{CAPTURA-AGUA-14.png}{Soporte de Agua con Subtipo instalación Agua visible en la paleta y en la confirmación previa a insertar}

\subsubsection{Tipos de soporte más representativos en Agua}

En la configuración actual, los casos más representativos en Agua son:

\begin{itemize}
\item \texttt{Omega}, cuando se trabaja con la lógica de soporte \texttt{Varifix}.
\item \texttt{Zeta}, cuando se necesita un soporte lineal o un soporte con desarrollo vertical.
\item \texttt{Cinta}, cuando se resuelve el soporte mediante cinta.
\end{itemize}

La diferencia en Agua no depende solo de la forma. También depende de cómo quede clasificado el soporte dentro de la instalación.

\subsubsection{Cotas que cambian el resultado en Agua}

En Agua, las cotas no son un detalle menor, porque cambian directamente el tipo de resultado que se obtiene:

\begin{itemize}
\item en \texttt{Omega} asociado a Agua, \texttt{Cota A} y \texttt{Cota B} deben ser mayores que \texttt{0};
\item en \texttt{Zeta}, \texttt{Cota B} debe ser mayor que \texttt{0};
\item en \texttt{Zeta}, si \texttt{Cota A = 0 mm}, el soporte se resuelve como variante lineal \texttt{0 mm};
\item en \texttt{Zeta}, si \texttt{Cota A} es mayor que \texttt{0 mm}, el soporte pasa a una variante con desarrollo vertical;
\item en \texttt{Cinta}, la medida decisiva es \texttt{Cota B}, que también debe estar informada.
\end{itemize}

Esto conviene revisarlo siempre en previsualización, porque un mismo tipo de soporte puede dar un resultado muy distinto solo por cambiar las cotas.

\manualimage{CAPTURA-AGUA-15.png}{Comparación en Agua entre un soporte Zeta con Cota A igual a 0 mm y otro con Cota A mayor que 0 mm}

\subsubsection{Layers e identificación de soporte en Agua}

Cuando el soporte queda correctamente asociado a Agua, la clasificación habitual de salida es:

\begin{itemize}
\item layer \texttt{IS\_SUPORTS\_AIGUA};
\item o \texttt{IS\_SUPORTS\_AGUA}, si el proyecto usa esa nomenclatura;
\item denominación de soporte \texttt{SUPORT AIGUA}.
\end{itemize}

En cambio, el \texttt{Zeta 0 mm} mantiene su condición de soporte lineal y conserva la lógica de layer lineal general, en lugar de pasar al layer específico de Agua.

Esto es importante porque permite distinguir entre:

\begin{itemize}
\item soportes de Agua ya clasificados como parte de esa instalación;
\item y soportes lineales genéricos que no deben interpretarse igual.
\end{itemize}

\manualimage{CAPTURA-AGUA-16.png}{Propiedades finales de un soporte de Agua mostrando su layer e identificación}

\subsubsection{Atributos de soporte en Agua}

En Agua, el \texttt{Atributo de soporte} no solo sirve para etiquetar visualmente el soporte. También ayuda a dejarlo correctamente identificado dentro de la instalación.

Si el flujo del proyecto utiliza organización por atributo padre, conviene revisar este valor antes de finalizar, porque puede influir en cómo se organiza el soporte dentro de la instalación.

Por eso, cuando los soportes formen parte de una misma sectorización o deban seguir la misma lógica que la instalación principal, es recomendable aplicar ese atributo antes de cerrar el trabajo.

\subsubsection{Qué conviene revisar antes de acumular o finalizar soportes en Agua}

\begin{itemize}
\item que \texttt{Subtipo instalación} sea realmente \texttt{Agua};
\item que un \texttt{Omega} no se haya quedado clasificado como otra instalación que use la misma familia de soporte;
\item que \texttt{Cota A} y \texttt{Cota B} correspondan al tipo de soporte elegido;
\item que un \texttt{Zeta 0 mm} no se haya usado por error cuando se necesitaba un soporte con desarrollo vertical;
\item que el layer y la identificación final del soporte sean los esperados;
\item y que el \texttt{Atributo de soporte} esté informado si el proyecto necesita clasificación por atributo.
\end{itemize}

\subsection{Puntos a revisar con especial atención en Agua}

Antes de finalizar una instalación de Agua conviene revisar especialmente:

\begin{itemize}
\item que el sistema seleccionado sea realmente \texttt{Polietilè}, \texttt{Multicapa} o \texttt{Armaflex}, según el caso;
\item que la distribución sea la correcta y, si procede, que \texttt{Cara (EN)} también esté bien definida;
\item que después de un cambio de diámetro se hayan regenerado correctamente codos, manguitos reductores y Tes, especialmente si la reducción aparece justo después de un codo;
\item que los giros del recorrido respondan a la lógica de \texttt{90 grados};
\item que los avisos de longitud mínima se hayan resuelto conscientemente;
\item que los layers de Agua se correspondan con el criterio de trabajo que se necesita en ese plano;
\item que el atributo padre se haya aplicado realmente a los elementos que deben compartir referencia;
\item que los soportes de Agua hayan quedado asociados al subtipo correcto, con sus cotas y layers bien resueltos;
\item y que los recorridos largos no hayan generado divisiones automáticas en puntos no deseados.
\end{itemize}

\section{Próximas unidades por instalación}

Este mismo esquema puede repetirse dentro del manual para el resto de instalaciones, por ejemplo:

\begin{itemize}
\item Saneamiento;
\item Ventilación;
\item Electricidad.
\end{itemize}

La idea es que cada una tenga su propia unidad específica dentro del mismo manual, manteniendo una estructura común para que la consulta sea sencilla.

\end{document}
